\chapter{Conclusion and Future Work}\label{chapter:conclusion}

\section{Conclusion}
This thesis studied organ-aware radiology report generation from 3D chest CT with a unified vision-language pipeline. The core objective was to reduce the gap between rich volumetric information and decoder-side token constraints while preserving clinically meaningful report content.

The implemented framework combined organ-level supervision via report decomposition, organ-conditioned visual token extraction, and decoder adaptation under a consistent evaluation protocol. The evidence now includes systematic MedGemma ablations and completed decoder-family rows for MedGemma, fVLM-aligned BERT-base, GPT-2, and BioBART under the same evaluation pass.

The main empirical conclusion is concise: in this setting, the combined choice of token budget and stable visual-language interfacing is more decisive than increasing architectural complexity. Within MedGemma ablations, the strongest model variant used a single-token visual interface with explicit visual-prefix structuring and delivered the most consistent improvements on the primary global metrics used in this thesis. Cross-decoder analysis also shows that stronger aggregate lexical-semantic metrics can coexist with higher template reuse, so deployment-oriented interpretation must include diversity and structure diagnostics in addition to metric averages.

Overall, the thesis contributes:
\begin{itemize}
    \item a reproducible end-to-end organ-aware 3D report generation setup,
    \item code-grounded ablation evidence on interface, architecture, and training strategy,
    \item a standardized evaluation workflow supporting multi-metric interpretation,
    \item a protocol-matched cross-decoder comparison with completed MedGemma, fVLM-aligned BERT-base, GPT-2, and BioBART rows.
\end{itemize}

The work therefore advances both method (how to structure the task) and practice (how to implement and evaluate it reliably).

\section{Take-Home Messages}
The thesis can be reduced to four practical takeaways:
\begin{enumerate}
    \item Organ-aware decomposition is not only a data-format choice; it is a training signal that makes failure analysis and controllable generation feasible.
    \item In 3D CT report generation, token budget and interface design can matter more than adding architectural complexity.
    \item Multi-metric evaluation is mandatory; single-metric optimization is too fragile for clinical text generation claims.
    \item Reproducible artifact discipline (\texttt{generated\_reports*.csv} + \texttt{metrics\_final.csv}) is essential for trustworthy model comparison.
\end{enumerate}

These messages are intended to be directly reusable by future work in volumetric medical vision-language systems, independent of decoder brand or training infrastructure.

\section{Future Work}
Future work is prioritized in three tracks.

\subsection{Track 1: Strengthen the Current Thesis Axis}
\begin{enumerate}
    \item Add multi-seed stability checks for the decoder-family table under the same standardized metric pass.
    \item Insert finalized cross-attention case-study figures to complete the quantitative + qualitative evidence pair.
    \item Add compact significance reporting (e.g., bootstrap confidence intervals) for key global metrics.
\end{enumerate}
This track is short-horizon and can be completed without redesigning the model stack. Its purpose is to close the evidence loop of the current thesis version.

\subsection{Track 2: Improve Model Robustness and Efficiency}
\begin{enumerate}
    \item Run controlled token-budget sweeps (1/2/4/8 tokens per organ) under matched compute.
    \item Compare mask downsampling strategies (area interpolation vs max-pooling) for mask-to-token alignment.
    \item Ablate bridge stabilization choices (with/without variance matching and visual positional embeddings).
    \item Compare loss composition (pure LM vs LM+alignment) on branches where alignment is implemented.
    \item Compare sampling policies (uniform vs organ-weighted vs hard-example weighted).
    \item Explore organ-adaptive token allocation and uncertainty-aware handling for sparse organs.
\end{enumerate}
This track targets the main technical bottlenecks identified by the ablations: token efficiency, interface stability, and long-tail organ robustness.

\subsection{Track 3: Strengthen Clinical Validity}
\begin{enumerate}
    \item Add human evaluation with radiologist-centered error taxonomy.
    \item Expand evaluation to temporal/longitudinal scenarios and clinically realistic reading workflows.
    \item Introduce safety-oriented post-generation checks for contradiction and hallucination control.
\end{enumerate}
This track is required for translational credibility. It shifts the system from benchmark-oriented progress to clinically meaningful utility and safety assessment.

\section{Final Remark}
The thesis shows that organ-aware 3D report generation is tractable when anatomical structure is treated as a first-class modeling constraint. With the current framework and completed decoder comparisons, the project is technically mature enough to move toward stronger clinical validation.
